\cotefile{156-7}
\scanpage{58}
\mspage{57}

Notons qu'on \deleted{\unsure{donc|doit}} doit avoir
\[
\begin{aligned}
&\deleted{\operatorname{omb}(X)\cap\operatorname{omb}(Y)\ne\varnothing}\\
&\operatorname{omb}(X)=\operatorname{omb}(Y)\qquad (=|\varphi(X)|=|\varphi(Y)|)
\end{aligned}
\]
et qu'a priori il \unsure{n'est} pas exclu que $L=\varnothing$.

\uline{Exemple} $X$ espace topologique. On prend \deleted{$\unsure{M}$}
\[
M=\text{simplexes \uline{affines} dans }X
\]
avec la relation $\leq$ habituelle. \deleted{$\ll$ \unsure{la relation}} (mais en respectant les structures affines), $\ll$ à l'\unsure{avenant} (sans exiger que $|\sigma|\cap|\sigma'|$ soit vide ou un sous-simplexe commun), et $\unsure{\gtrless|\lessgtr|\gtreqless|\lesseqgtr}$ itou ($\unsure{\simeq|\cong}$ \unsure{presque}), d'où $\unsure{\mathscr{F}|\mathscr{E}}$, satisfaisant toutes les conditions\inserted{ de (ii)}. Mais considérons, sur une partie de $X$ homéomorphe à $[0,1]$, deux structures \unsure{affines} \deleted{\unsure{deux}} différentes, d'où $X,Y\in M$ avec \inserted{\unsure{isom.}} $L=\partial X=\partial Y$, $\varphi(X)=\varphi(Y)$, mais $X\ne Y$. Donc ici
\[
\varphi(X)\gtrless\varphi(Y)\quad\text{et}\quad\varphi(X)\ll\varphi(Y),
\]
sans avoir ni
\[
X\gtrless Y\quad\text{ni}\quad X\ll Y.
\]

% margin note written sideways; slanted upward along the left edge
\margin{Plus précisément\newline $(i)\Rightarrow(ii)+\varphi\text{ injectif}\Rightarrow(i\,a,c)^\circ$}

% A vertical marginal stroke brackets the corollary.
\uline{Cor} La condition (i) \deleted{\unsure{équivaut à}} \inserted{implique} (ii) + régularité de $\varphi$. \deleted{Il suffit \unsure{de} d'avoir l'injectivité de $\varphi|M$,} \inserted{Mais l'inverse n'est} pas vrai.

Bien sûr, i implique l'injectivité de $\varphi$, reste à prouver que $(ii)+\varphi|M\ \text{inj.}\Rightarrow b,c$ \unsure{(}ça a été prouvé déjà).

Prouvons c), plutôt, continuons. On a vu (utilisant l'injectivité de $\varphi$) que

\scanpage{59}
\mspage{58}

$\varphi(L)$ est la plus grande sous-figure commune de $\varphi(F),\varphi(G)$. L'hyp. $\varphi(F)\unsure{\gtrless|\lessgtr|\gtreqless|\lesseqgtr}\varphi(G)$ équivaut à
\[
\begin{aligned}
&|\varphi(F)|\cap|\varphi(G)|=|\varphi(L)|\qquad\text{i.e.}\\
&\operatorname{omb}(F)\cap\operatorname{omb}(G)=\operatorname{omb}(L),
\end{aligned}
\]
mais c'est (en vertu de \unsure{AtCL (ap a)|AtCL (op a)}) la condition pour que $F\gtrless G$. OK.

Prouvons b), i.e. soient $F,G$ avec $\varphi(F)\ll\varphi(G)$, prouvons $F\ll G$. \deleted{\unsure{Il}} Utilisant a) déjà prouvé, on est ramené au cas $X,Y\in M$.

Mais c'est encore faux!

\uline{Ex} \deleted{\deleted{$X$ espace affine sur $\unsure{\mathbb{R}|R}$. On prend $M\unsure{=|-}$ simplexes affines, $\leq$ \unsure{pour} sous-simplexe,} \inserted{On reprend $\unsure{M}$ comme dans l'exemple précédent,} mais en renforçant cette fois la relation $\ll$, pour la rendre conforme à la version ordinaire (``\unsure{simpliciale}'' ou ``polyédrale''\unsure{.})}

Soit $L$ un ens., considérons l'ens $\unsure{\mathbb{F}|\mathrm{F}}\mathrm{ig}\,\mathrm{pol}(L)$ des figures \uline{polyédrales} dans $L$. On y prend $\leq$ et $\ll$ comme d'habitude\inserted{, dans le cas polyédral}, étant entendu que $\ll$ signifie $\overset{\mathrm{pol}}{\ll}$, tandis qu'on prend la compatibilité $\gtrless$ \uline{\unsure{ens}} (\unsure{que|par} \unsure{sur|sous} ensembles \unsure{en|ou} polyèdres). (D'avoir renforcé le $\underset{\mathrm{ens}}{\ll}$ en $\underset{\mathrm{pol}}{\ll}$ ne modifie pas la validité des conditions $\unsure{\mathrm{Prat}|\mathrm{Prot}}\,1'-\mathrm{Prat}\,6'$.) Considérons deux \deleted{figures}\inserted{str. de} polyèdres (p. ex. de \deleted{\unsure{qui sont}} type $\Delta_1,\unsure{\Delta_n|\Delta_1}$) $X,Y$, avec $X\underset{\mathrm{ens}}{\ll}Y$ et non $X\underset{\mathrm{pol}}{\ll}Y$, on a \unsure{le|là} contre-exemple.

\scanpage{60}
\mspage{59}

\uline{\unsure{Def|Déf}} Un atelier est dit ``ensembliste \unsure{spacieux|spécieux}'' s'il satisfait la condition (i). Idem pour ``ensembliste polyédral'', en remplaçant $\mathrm{Fig}(L)$ par $\mathrm{Fig}\,\mathrm{pol}(L)$. On a les \deleted{conditions} notions correspondantes pour les préateliers.

% margin note written sideways; slanted upward
\margin{NB. Il faut voir les figures dans un tel atelier comme étant des figures ensemblistes dans $L$, munies de \unsure{structures suppl\textsuperscript{taires}|structures supplémentaires} (\unsure{tout} comme les \unsure{2-figures|S-figures} en $L$, exemple plus haut).}

Un atelier est dit ``quasi-ensembliste \unsure{spacieux|spécieux}'' s'il satisfait (ii). Notions correspondantes pour le cas ``quasi-ensembliste polyédral'', ou le cas des préateliers.

Les préateliers quasi-ensemblistes (indifféremment \unsure{spacieux|spécieux} ou polyédraux) correspondent aux $(M,\leq,\ll)$ satisfaisant les conditions suivantes (à déduire de $\unsure{\mathrm{Prat}|\mathrm{Prot}}\,1'-\mathrm{Prat}\,6'$) :

\uline{Prat quas} 1
\[
X\leq Y\ \Rightarrow\ X\ll Y
\]

% slanted margin labels, with arrows toward the numbered conditions
\margin{densité\newline divisibilité $\Longrightarrow$}
\uline{Prat quas} 2
\[
\forall X\in M,\quad\operatorname{omb}(X)^\circ\ne\varnothing\quad\text{i.e.}\quad\exists x\in L,\quad x\overset{\circ}{\ll}X
\]

\margin{\unsure{Aut. minimale} $\longrightarrow$}
\uline{Prat quas} 3
\[
\forall X\in M,\ x\in L,\ x\ll X,\quad\exists X'\in M,\quad x\overset{\circ}{\ll}X'\leq X
\]

\margin{raffinement induit $\longrightarrow$}
\uline{Prat quas} 4 Soit
\diagram{\begin{tikzcd}[ampersand replacement=\&,column sep=small,row sep=small]
Y \arrow[r,phantom,"\ll" description] \& X \arrow[d,phantom,"\unsure{\vee|\geq}" description] \\
\& X'
\end{tikzcd}}
\deleted{Alors} et $Y\deleted{\unsure{X'}}_{X'}=\widetilde Y\cap X'$.

Alors
\[
\operatorname{omb}(Y)\cap\operatorname{omb}(X')=\operatorname{omb}(\unsure{Y'|Y_{X'}})\left(\overset{\mathrm{def}}{=}\bigcup_{Z\in Y_{X'}}\operatorname{omb}(Z)\right),
\]
i.e. si
\[
x\in L,\quad x\ll Y,X',\qquad\exists Y'\text{ avec }Y'\leq Y,\ Y'\ll X',\ x\ll Y'.
\]
De plus, si $Y\leq X$, OPS $Y'\leq X'$ (i.e. on peut remplacer $Y_{X'}$ par $\widetilde Y\cap\widetilde X'$ : -- les deux \unsure{seraient|seront} d'ailleurs égaux).

Il n'y a plus lieu en effet d'introduire $\parallel_L$, vu qu'on la connaît : \deleted{\unsure{c.à.d.}} $(x\parallel y)\Longleftrightarrow(x\ne y)$.

\scanpage{61}
\mspage{60}

Pratiquement, quand on ne travaille pas avec des ``lieux'' proprement dits (et je suis le premier, apparemment, à y avoir songé, aux lieux\ldots) mais avec des points, et dans l'optique des espaces ``infiniment divisibles'' (et avec \unsure{bouzy|Bouzy|Bourb.}, je suis le premier à songer à travailler \unsure{éo!|eô!|eo!} avec des espaces qui ne le soient pas), il semblerait que \deleted{\unsure{ça soit bon}} cette structure \unsure{satisfasse|satisfera} à tous les besoins (avec, bien sûr, des axiomes \uline{en plus}, suivant les besoins). Cette structure a l'avantage d'être \uline{commune} aux contextes ``\unsure{spacieux|spécieux}'' et ``polyédral'', et se recommande de plus par la simplicité étonnante des \uline{données} (seulement $\leq$ et $\ll$ dans $M$ \unsure{lui-même} dans le cas des ateliers $F$), et des axiomes, réduits ici au nb de quatre (à peine plus compliqué qu'une topologie!). Quand on travaille avec des \deleted{\unsure{Quand}} ateliers à figures finies\inserted{ (cas le plus important)}, ceux-ci sont définis déjà par le préatelier correspondant.

% margin note written sideways; slanted upward, crossing the left edge of the paragraph
\margin{Ça mériterait peut-être un nom plus lapidaire que ``Préatelier quasi-ensembliste'', p.ex. ``\unsure{Stratop}'', ou (avec moins d'emphase) ``Stratologie'' (à moins d'utiliser ``Préatelier'' sans plus). L'usage dira!}

Si on veut travailler avec des ateliers qu. ens. (\unsure{spacieux|spécieux} resp. polyédraux) plus généraux, il faut introduire sur $M$ la notion de \deleted{\unsure{famille}} partie (fermée) ``admissible'' par des conditions convenables de ``locale finitude''.

\scanpage{62}
\mspage{61}

\uline{Notion de support.}

Soit $\unsure{\mathscr{F}|\mathcal{F}}$ un atelier. Soit $\Sigma_{\mathscr{F}}\subset\unsure{\mathcal{P}|P}(\mathscr{F})$ l'ens. des parties de $\mathscr{F}$ qui sont fermées \deleted{à la fois} pour \deleted{$\leq$ et} $\ll$\inserted{ (donc aussi pour $\leq$)}, et stables par Sup pour $\leq$ (Sup finis ou infinis). Soit $\Sigma_M$ l'ens des parties de $M$ fermées pour $\ll$. Alors on a deux bijections inverses l'une de l'autre

\noindent $(\unsure{\alpha|1|*})$
\diagram{\begin{tikzcd}[ampersand replacement=\&,column sep=large,row sep=small]
\Phi \arrow[r,mapsto] \& \Phi\cap M\,\deleted{\unsure{\subset|\in}} \\
\Sigma_{\mathscr{F}} \arrow[r,shift left=1.1ex] \& \Sigma_{\deleted{\unsure{\mathscr{F}}}M} \arrow[l,shift left=1.1ex] \\
\Phi_S=\{F\in\mathscr{F}\mid\widetilde F\subset S\} \& \deleted{\Phi_{\unsure{S|c}}}S \arrow[l,mapsto]
\end{tikzcd}}
\margin{\uline{Iso. d'ens.}\newline\uline{ordonnés}}

moyennant quoi on va les identifier, trouvant un ens. $\Sigma$.

Soit $A\subset\mathscr{F}$ une partie quelconque. On définit son \uline{cosupport} par
\[
\operatorname{cosupp}A=\{G\in\mathscr{F}\mid\forall F\in A,\ F\parallel G\}\subset\mathscr{F}
\]
On voit, grâce à $\unsure{\mathrm{AtD}_0|\mathrm{AtD}_o}$, que
\[
\operatorname{cosupp}A\in\Sigma_{\mathscr{F}}
\]
\deleted{et il \unsure{est}}

% margin note written sideways; slanted upward
\margin{On écrit $F\parallel A$ au lieu de $F\in\operatorname{cosupp}A$.}

On \deleted{\unsure{écrit}} pose
\[
\operatorname{supp}A\inserted{\ (\text{ou }\overline A)}=\operatorname{cosupp}(\operatorname{cosupp}A),
\]
\margin{Ne pas confondre $\overline A$ avec l'adh. de $A$ pour la top. de $\mathscr{F}$ associée soit à $\leq$, soit à $\ll$!}

% margin note written sideways; slanted upward
\margin{On écrit $A\parallel B$ pour les relations équivalentes\newline
$(i)\ A\subset\operatorname{cosupp}B$\newline
$(ii)\ B\subset\operatorname{cosupp}A$\newline
\deleted{\unsure{(iii)}}\unsure{(ii')} $\forall X\in A,\ X\parallel B$\newline
$(iii)\ \forall Y\in B,\ Y\parallel A$\newline
$(iii')\ \forall X\in A,Y\in B,\ X\parallel Y$.\newline
\unsure{Ser.} équivalent aussi (\unsure{cf|ce} dans le cas \unsure{ensembliste}) $\operatorname{supp}A\cap\operatorname{supp}B=\varnothing$\deleted{$\unsure{\mathscr{F}}$} (ça l'implique en tous cas).}

Il est formel (en termes de la seule donnée d'une relation $\parallel$ dans l'ens abstrait $\mathscr{F}$ -- \unsure{ça} marcherait aussi pour une relation quelc. entre deux ens)

\scanpage{63}
\mspage{62}

\[
\deleted{\unsure{A\subset\overline A},}
\]
\[
A\subset B\ \Longrightarrow\ \operatorname{cosupp}A\supset\operatorname{cosupp}B
\tag*{(i)}
\]
\[
A\subset\overline A,\quad\text{et }A\subset B\ \Longrightarrow\ \overline A\subset\overline B,
\tag*{(ii)}
\]
\margin{donc $\displaystyle\bigcup_i\overline A_{\unsure{j|i}}\subset\overline{\bigcup A_i}$ pour toute famille d'ens. (mais on n'a pas en général égalité)}
\deleted{\unsure{(iii)}}
\[
\overline{\overline A}=\overline A
\tag*{(iii)}
\]

(iv) Conditions équivalentes sur $A\in\unsure{\mathscr{F}|\mathcal{F}}$.
\[
\begin{aligned}
\text{a)}\quad&\exists B\text{ avec }A=\operatorname{cosupp}B\\
\text{b)}\quad&A=\overline A
\end{aligned}
\]
\deleted{\unsure{c)}}

On dira alors que $A$ est une partie \uline{supp-fermée} de $\mathscr{F}$, \deleted{\unsure{et}} ou ``\inserted{\uline{un}} \uline{support}''. \deleted{\unsure{Alors}} On a alors $A\in\Sigma_{\mathscr{F}}$.

\deleted{(v) $\forall A\subset\mathscr{F}$, $\overline A$ est le plus}

Soit $\deleted{\Sigma^\circ}\ \Sigma_{\mathscr{F}}^\circ=\deleted{\Sigma^\circ}$ la partie de $\Sigma_{\mathscr{F}}$ formée des \deleted{parties \inserted{supp} fermées de $\mathscr{F}$. supports de $\mathscr{F}$.} $\mathscr{F}$-supports. On a alors

(v) Si $A\subset\mathscr{F}$, $\overline A$ est le plus petite \deleted{partie} \deleted{$\mathscr{F}$-fermée} $\mathscr{F}$-support contenant $A$.

(vi) L'application
\diagram{\begin{tikzcd}[ampersand replacement=\&,column sep=large]
A \arrow[r] \& \operatorname{cosupp}A \& \unsure{\mathcal{P}|P}(\mathscr{F}) \arrow[r] \& \Sigma_{\mathscr{F}}^\circ
\end{tikzcd}}
\deleted{\unsure{est une application} de $\unsure{\mathcal{P}(}$}
induit un \inserted{anti} Iso. d'ens. ordonnés
\diagram{\begin{tikzcd}[ampersand replacement=\&]
\Sigma_{\mathscr{F}}^\circ \arrow[r] \& \Sigma_{\mathscr{F}}^\circ
\end{tikzcd}}
qui est \uline{involutif} (exprimant le fait que la relation $\parallel$ est symétrique). On pose aussi, si $A$ est un support
\[
\deleted{\unsure{C}}\unsure{\complement|\mathcal{C}}A=\operatorname{cosupp}A,
\]

\scanpage{64}
\mspage{63}

on l'appelle le \uline{support complémentaire} du support donné $A$.

(vii) $\Sigma_{\unsure{\mathscr{F}|\mathcal{F}}}^\circ$ a un plus petit élément $\varnothing$ (à savoir la partie \deleted{vide} de $\mathscr{F}$\inserted{ réduite à $\varnothing_{\mathscr{F}}$*}), et un plus grand élément (à savoir $\mathscr{F}$), qui sont complémentaires l'un de l'autre.

% margin note written sideways; slanted upward
\margin{* Ceci provient du fait que $\varnothing_{\mathscr{F}}$ est l'unique élément de $\mathscr{F}$ qui soit $\parallel$ de tout élément de $\mathscr{F}$ (ou \unsure{encore} disjoint \unsure{d'elle-même|de lui-même}). On le notera, par abus de langage, $\varnothing$ (cf plus bas).\newline $\varnothing_\Sigma$}

(viii) \inserted{Plus gén.} Dans $\Sigma_{\mathscr{F}}^\circ$, il y a des Sup et des Inf quelconques, échangés par l'application $\unsure{\complement|\mathcal{C}}$. Le \deleted{\unsure{produit}} inf s'explicite par
\[
\underset{i}{\operatorname{Inf}}^{\Sigma_{\mathscr{F}}^\circ}A_i=\bigcap A_i
\]
et le Sup par
\[
\underset{i}{\operatorname{Sup}}^{\Sigma_{\mathscr{F}}^\circ}A_i=\overline{\bigcup A_i}\quad\Rightarrow\ (\text{ou }\operatorname{supp}(\bigcup A_i))
\]

\unsure{(iix)|(ix)} $\forall A\in\Sigma_{\mathscr{F}}^\circ$, on a
\[
A\cap\complement A=\varnothing_\Sigma
\]
\deleted{(En effet si $\unsure{F\in A}$ \unsure{de} $\varnothing_{\mathscr{F}}$}
(Exprimant le fait que le seul élément de $\mathscr{F}$ \deleted{\unsure{qui}} tel que $F\parallel F$ est $\varnothing_{\mathscr{F}}$, lequel est $\parallel$ de $\mathscr{F}$ tout entier\ldots)

($\unsure{x}$) \uline{Cor}
\[
A\subset B\ \Longrightarrow\ A\cap\complement B=\varnothing_\Sigma
\]

\uline{NB} Il n'est pas clair que l'implication inverse soit vraie. \inserted{$\Rightarrow$ \unsure{l'est}} Il n'est pas clair \deleted{que} (il est \deleted{parfois} faux) que le Sup\inserted{ de $\Sigma_{\mathscr{F}}^\circ$} \deleted{soit} s'interprète comme la réunion dans $\mathscr{F}$ -- la réunion\inserted{ (en tant que parties de $\mathscr{F}$)} de deux supports, n'est pas forcément un support : cela exprime le fait\inserted{ intuitif} qu'une ``figure''\inserted{ (dont le support est)} (\unsure{contenu}\deleted{\unsure{e}}) dans la réunion des supports de deux figures n'est pas forcément contenue dans\inserted{ (le support de)} l'une d'elles -- même dans le cas des multistrates, et $\unsure{=|\overline m}$ dans le cas des

\scanpage{65}
\mspage{64}

lieux (quand on n'est pas dans un atelier ensembliste tout au moins\ldots). Il n'est pas clair non plus, en général, que l'on ait, dans $\Sigma_{\unsure{\mathscr{F}|\mathcal{F}}}^\circ$,
\[
P\unsure{\wedge|\cap}\deleted{\unsure{\operatorname{Sup}_\Sigma|\operatorname{Sup}_{\mathscr{F}}}}\bigvee_i Q_i=\deleted{\unsure{\mathscr{F}|\Sigma|\operatorname{Sup}}}\bigvee_i(P\wedge Q_i)
\]
(distributivité de l'intersection des supports, par rapport aux ``sommes'' même finies\ldots).
Pourtant, cela a tendance à être vrai, de $\unsure{=|\overline m}$ que l'implication inverse de (x), dans les \deleted{\unsure{lieux où j'ai été}}\inserted{ateliers auxquels j'ai songé} jusqu'à présent, cf plus bas.

Tout \unsure{ceci|ce soit} se répéterait dans le cadre des \unsure{Stratos} $(M,\leq,\ll,\unsure{\gtrless|\lessgtr|\gtreqless|\lesseqgtr})$, \deleted{et} en termes de la seule relation $\parallel_M$ \deleted{et} de disjonction dans $M$. Et l'ens. $\Sigma_M^\circ$ des ``$M$-supports'' est en bij. canonique avec son homologue
\diagram{\begin{tikzcd}[ampersand replacement=\&]
\Sigma_M^\circ \arrow[r,shift left] \& \Sigma_{\mathscr{F}}^\circ \arrow[l,shift left]
\end{tikzcd}}
\margin{iso. d'ens. ord.}
(induit par les bijections de la p. 61)

% margin note written sideways; slanted upward
\margin{Si $A\subset\mathscr{F}$ est fermée pour $\leq$, on a\newline
$\operatorname{cosupp}_{\mathscr{F}}(A)\cap\unsure{L|M}=\operatorname{cosupp}_L(A\cap L)$,\newline
et pour $F\in\mathscr{F}$, on a\newline
$\operatorname{cosupp}_{\mathscr{F}}(F)\cap L=\operatorname{cosupp}_L\unsure{C^\circ(F)|\mathcal{C}^\circ(F)|(\mathring{\widetilde F})}$.}

commutant à l'opération ``complémentaire''.
Cette fois, le plus petit élément de $\Sigma_M^\circ$ est bien la partie \uline{vide} de $M$, et son plus grand élément est $M$.
